\section{Problem Formulation}
\label{sec:problem}

We formalize individual-level silicon sampling as a conditional choice prediction problem. Let $\mathcal{D}=\{(r,q,o,y)\}$ be a survey dataset where $r$ indexes a respondent, $q$ is a question stem, $o=(o_1,\dots,o_K)$ is the ordered option set with $K$ alternatives, and $y\in\{1,\dots,K\}$ is the respondent's true choice. Each respondent $r$ is associated with a demographic profile $D_r$ (e.g., age, gender, education, region, income bracket) and a set of prior answer records $H_r=\{(q_i,a_i,o_i)\}_{i=1}^{n_r}$, where each record retains the full option set $o_i$ alongside the chosen answer $a_i$. Given a target item $(r,q^\ast,o^\ast)$ unseen at training time, the goal is to predict $y^\ast$ accurately, with calibrated uncertainty, in a manner that generalizes across items, domains, and demographic intersections.

\paragraph{The structural failure of direct generation.} The dominant paradigm reduces this to a sequence-to-token mapping $f_{\text{LLM}}(D_r,H_r,q^\ast)\to y^\ast$, in which the LLM is prompted to emit a single answer token. We identify three structural defects of this reduction. First, the answer token is a point estimate that discards the geometry of the respondent's latent reaction distribution over the option simplex; two respondents whose reaction distributions differ sharply in their second-ranked option but agree on the top option are indistinguishable after tokenization. Second, sampled tokens inherit the LLM's response bias and are poorly calibrated against the empirical choice distribution of the target instrument~\citep{geng2024survey}. Third, prompt-to-token pipelines couple the model's world knowledge to the specific instrument, so that prompt engineering tuned on one item set does not transfer to unseen items, other domains, or unseen demographic combinations~\citep{ku2026silicon}.

\paragraph{Representation-based reformulation.} We propose instead to factor the prediction into two stages, replacing the direct mapping $f_{\text{LLM}}$ with
\begin{equation}
\label{eq:factorization}
\hat{y}^\ast = h_\phi\!\left(\, g_\theta(D_r,H_r,q^\ast,o^\ast),\; o^\ast \,\right),
\end{equation}
where $g_\theta$ is an LLM-based representation generator producing a latent vector, and $h_\phi$ is a lightweight, instrument-aware decoder. This factorization mirrors the move from predictor to representation generator demonstrated in spatio-temporal learning by \citet{he2025geolocation}, who showed that extracting an LLM-derived geolocation vector and feeding it to a small downstream model yields a generic enhancer that transfers across cities, far outperforming direct LLM forecasting. Our transposition to the survey domain rests on the same insight: the LLM carries rich information about how a respondent would react, but committing to a discrete answer is a lossy readout; reading out a vector instead preserves the structure needed for downstream calibration and transfer.

\paragraph{Generative vs.\ discriminative encoders.} The choice of $g_\theta$ is consequential. A discriminative contrastive encoder such as LLM2Vec~\citep{behnamghader2024llmvec} maps inputs to a representation space optimized for retrieval, discarding the LLM's output distribution semantics. A generative embedder such as LLM2Vec-Gen~\citep{behnamghader2026llmvecgen} preserves those semantics by training the embedder to retain the model's generative output. Because the object we wish to represent---a latent disposition to answer a question in a particular way---is intrinsically tied to the LLM's generative distribution over the option vocabulary, we argue that a generative embedder is the correct inductive bias. We contrast this against raw hidden states (which are geometrically anisotropic and unaligned to choice) and against generic text vectors~\citep{wang2024multilingual} (which are not option-aware).

\paragraph{Two representation regimes.} Equation~\ref{eq:factorization} admits two regimes that we study in tandem. In the \emph{query-conditioned} regime, $g_\theta$ consumes the target item $(q^\ast,o^\ast)$ together with $(D_r,H_r)$, producing a per-item \emph{ResponseVec} $\mathbf{v}^\ast_r$. In the \emph{reusable} regime, $g_\theta$ consumes only $(D_r,H_r)$, producing a per-person \emph{RespondentVec} $\mathbf{v}_r$ that is computed once and reused across all items; the target item enters only at the decoder. The two regimes expose the central tension we study: query conditioning maximizes task adaptivity but incurs a per-item LLM forward pass, whereas reuse amortizes compute but sacrifices item specificity. Section~\ref{sec:tradeoff} formalizes this boundary.
